\documentclass[10pt,a4paper]{article}

% ============================================================
% PACKAGES
% ============================================================
\usepackage[top=0.5in, bottom=0.45in, left=0.7in, right=0.7in]{geometry}
\usepackage{titlesec}
\usepackage{enumitem}
\usepackage{hyperref}
\usepackage{xcolor}
\usepackage[T1]{fontenc}
\usepackage{charter}
\usepackage{parskip}
\usepackage{needspace}
\usepackage{ragged2e}

% ============================================================
% COLORS (match raadh.me)
% ============================================================
\definecolor{heading}{HTML}{0B4F5C}
\definecolor{link}{HTML}{0B4F5C}
\definecolor{body}{HTML}{1A1A1A}
\definecolor{muted}{HTML}{5A6470}

% ============================================================
% HYPERLINKS
% ============================================================
\hypersetup{
    colorlinks=true,
    linkcolor=link,
    urlcolor=link,
    pdfauthor={Rokib Al Dhin Raadh},
    pdftitle={Rokib Al Dhin Raadh - Cover Letter (Research Intern, Basis)},
    pdfsubject={AI researcher; founder and CEO of OXIEDO; inventor of ORMAS; DeepMath 2026; NeurIPS 2026 Workshop reviewer; Cohere Labs invited talk},
    pdfkeywords={Rokib Al Dhin Raadh, model evaluation, training observability, regression detection, failure modes, held-out evaluation, ablations, PyTorch, LLMs, ORMAS},
}

% ============================================================
% SECTION STYLE
% ============================================================
\titleformat{\section}
  {\normalsize\bfseries\scshape\color{heading}}
  {}
  {0em}
  {}
  [\vspace{-1pt}{\color{heading}\titlerule[0.4pt]}]
\titlespacing{\section}{0pt}{9pt}{5pt}

\pagestyle{empty}

% ============================================================
% CV ENTRY: date in the left gutter, body indented
% ============================================================
% A list (not a minipage) so long entries can break across pages.
\newcommand{\cvitem}[3]{%
  \begin{list}{}{%
    \setlength{\leftmargin}{2.35cm}%
    \setlength{\labelwidth}{2.1cm}%
    \setlength{\labelsep}{0.25cm}%
    \setlength{\itemindent}{0pt}%
    \setlength{\listparindent}{0pt}%
    \setlength{\rightmargin}{0pt}%
    \setlength{\topsep}{0pt}%
    \setlength{\partopsep}{0pt}%
    \setlength{\parsep}{0pt}%
    \setlength{\itemsep}{0pt}%
  }%
  \item[{\makebox[2.1cm][l]{\footnotesize\color{muted}#1}}]%
  #2
  #3
  \end{list}%
  \vspace{1.5pt}
}

% Tight bullet list
\newlist{bullets}{itemize}{1}
\setlist[bullets]{leftmargin=1em, itemsep=0pt, parsep=0pt, topsep=1pt, label=\textbullet}

% Numbered publication list
\newlist{pubs}{enumerate}{1}
\setlist[pubs]{leftmargin=2.35cm, labelwidth=2.1cm, labelsep=0.25cm, labelindent=0pt, align=left, itemsep=2pt, parsep=0pt, topsep=1pt, label={\footnotesize\color{muted}[\arabic*]}}

\setlength{\emergencystretch}{2.5em}
\hyphenpenalty=10000
\exhyphenpenalty=50
\begin{document}
\RaggedRight
\color{body}
\setlength{\parskip}{7pt}
{\large\bfseries\color{heading} Rokib Al Dhin Raadh}\\[1pt]
{\footnotesize \href{mailto:raadh@oxiedo.com}{raadh@oxiedo.com} \enspace$|$\enspace \href{https://www.raadh.me}{raadh.me} \enspace$|$\enspace \href{https://github.com/Oxiedo/ormas-monitor}{github.com/Oxiedo/ormas-monitor} \enspace$|$\enspace Dhaka, Bangladesh}

\vspace{6pt}
{\small Basis \enspace$\cdot$\enspace Research Intern, Summer 2027 \enspace$\cdot$\enspace Dynamical Systems + Machine Learning, or Formal Methods and Verification}

\vspace{4pt}
Dear Basis team,

Your rollback work for safer agents rests on one rule: before a state-changing call goes through, record how to undo it, and refuse any call you cannot undo. I reached the same idea from inside neural network training, and found that some changes cannot be undone at all.

I model training as a dynamical system in which a self-repair mechanism acts as a bounded disturbance on gradient descent. In my sole-author DeepMath 2026 paper (Ohio State, double-blind review), I proved local practical input-to-state stability for this system. I also proved a negative result: once a BatchNorm--ReLU channel dies, no per-channel-bounded repair can bring it back. That set of states is absorbing. Some damage during training is permanent, and you can certify it.

I then tried to break the result on real models. Turned into a check, it found near-zero channels in 19 of the 21 torchvision models that use BatchNorm and ReLU. In VGG-19-BN, 292 channels output exactly zero on all 10,000 ImageNetV2 photos. Removing them cut 18.6\% of the compute and changed none of the predictions. Removing 292 random live channels instead changed 499 of 500. I released this as \href{https://pypi.org/project/ormas-monitor/}{ormas-monitor}, which certifies dead channels and keeps a tamper-evident log of every correction made during training, so the evidence of how a model changed is kept as it evolves.

That is why both of your tracks fit me. In Dynamical Systems + Machine Learning, I want to learn the inference side I don't yet have: uncertainty quantification and system identification, the questions Dynestyx is built for. In Formal Methods and Verification, I'd like to turn my guarantees into machine-checkable ones instead of proofs on paper.

I've completed a full research cycle alone: theory, 383 seeded experiments on one GPU, a peer-reviewed paper, an open-source tool, and an invited talk at Cohere Labs where I presented corrections to my own results. What I haven't had is a team that works on foundations and on real problems. Basis is the rare place that does both.

To be direct about logistics: I'm 18 and finishing high school in Bangladesh in 2027, and I would need visa sponsorship. If summer 2027 isn't workable, I would be glad to be considered for 2028, or to contribute to your open-source work remotely in the meantime.

Thank you for reading.

Rokib Al Dhin Raadh
\end{document}
